\section{Discussion}
\label{sec:discussion}

\subsection{Limitations}
\label{sec:disc-limits}

The framework's reach is bounded by its reliance on a non-LLM verifier. The
agent learns from reward, and reward in our setting must come from a checker
that is external to the language model: a symbolic solver, a unit-test harness,
a proof kernel. Domains in which no such checker exists fall outside the
present scope. Open-ended deliberation tasks such as essay writing, persuasive
dialogue, or therapeutic conversation cannot be cast into this framework
without adding a learned signal back into the reward channel, and doing so
would surrender the very Goodhart-resistance property
(Section~\ref{sec:disc-verifier}) that motivates the design. We treat this as
a deliberate restriction rather than an oversight: the contribution is sharper
when the reward signal is unimpeachable, and the verifiable-math setting is
where the methodology can be cleanly evaluated.

A second restriction stems from the phase 1 decision to freeze the base
language model. If the latent capability required for a problem class is
absent from the frozen base, the deliberation loop cannot manufacture it. The
heads can re-weight the latent skills the model already possesses, schedule
verifier calls, and prune unprofitable branches; they cannot teach the
underlying network to compute things it has never been trained to compute.
Phase 2 admits unfreezing of the last one or two transformer layers via LoRA
\citep{hu2022lora}, but that is a separate intervention whose contribution
must be measured against the frozen baseline. Closely related is our use of a
fixed verb set. Problems whose solutions require verbs we did
not anticipate (literature search, tool invocation outside the verifier set,
mid-stream consultation of an external oracle) are out of reach until
migration to progressive widening with LLM-proposed actions
(Section~\ref{sec:method-actions}).

Two further limitations follow from the structural commitments of the
framework. The Markov bet that all decision-relevant information lives in the
structured state schema (Section~\ref{sec:method-state}) places a load-bearing
demand on the parser that lifts LLM-generated text into \texttt{Fact},
\texttt{Hypothesis}, and \texttt{VerifierCall} records. Lossy parsing surfaces
downstream as apparent $Q$-head failures: when a $Q$-value disagrees with
expectation, the parser is the first suspect, not the head. We accept the bet
because the alternative (raw token sequences as states) forecloses any
structural inductive bias and any tree reuse. Finally, the experimental plan
calls for multiple training seeds, but compute may force a tradeoff between
seed breadth and simulation budget at the largest configurations. A
single-seed result is not, by itself, evidence that the framework works; we
take this as a known threat to the eventual empirical story rather than a
defect of the framework as specified.

\subsection{Extensions}
\label{sec:disc-extensions}

The clearest near-term extension is progressive widening with LLM-proposed
actions, sketched as option C in Section~\ref{sec:method-actions}. Replacing
the categorical multi-head $Q$ with an $(s, a)$-embedding-based scorer allows
the action space to grow with the visit count and admits domains where the
useful verbs are themselves unknown at design time. The architectural change
is bounded: the adapter, serialization, and verifier-anchored reward survive
unchanged; only the $Q$-head's interface changes, from a fixed categorical
output to a scoring function over candidate $(s, a)$ pairs sampled by the LLM.
This is the pattern of AlphaProof \citep{deepmind2024alphaproof} and AlphaCode
\citep{li2022alphacode}, and our framework accepts it as a local change rather
than a redesign.

Several smaller extensions follow naturally. Canonicalizing states (topological
ordering of facts by derivation depth, SymPy normal forms on statements, sorted
subgoal identifiers) would let MCTS reuse subtrees across trajectories that
reach semantically equivalent contexts in different orders; deferred in phase 1
for debuggability, it is a multiplicative win in phase 2 for problems with rich
equivalence classes among intermediate states. If frozen-encoder approximate
$Q$-learning underfits, LoRA on the last one or two transformer layers
\citep{hu2022lora} is the natural next step, its contribution measured against
the frozen baseline. Per-channel $Q$ heads (reward decomposed into correctness,
safety, and instruction-following components, each with its own verifier)
generalize the single-channel correctness signal to domains where ``the right
answer'' is no longer a single Boolean; adding channels later is a head-shape
change, not a redesign of the loop.

Deployment-side extensions also fit cleanly. Persistent search trees that
survive across queries support anytime deliberation: a caller who asks the
system to ``think harder'' gets more simulations grafted onto the existing
tree, and rerooting on a committed child while reusing surviving subtrees
generalizes the same machinery to multi-turn deliberation. Online learning
from user-supplied verifier signals (a developer marking ``this proof has a
gap'') is feasible as long as the verifier remains a structural check rather
than a learned classifier; the risk that one user's signals dominate updates
is mitigated by adapter-tier separation, where per-user adapters absorb
idiosyncratic feedback, per-cohort adapters aggregate it, and the global
adapter accepts only changes that survive both tiers. The framework supplies
the place where these adapters plug in (the small adapter MLP of
Section~\ref{sec:method-architecture}); building out the tiers is engineering
the methodology does not need to specify.

\subsection{Why Verifier-Anchored Signals Are Different}
\label{sec:disc-verifier}

The asymmetry between verifier-anchored reward and learned-aesthetic reward
is the most consequential design choice in this paper, and we develop it in
full here because the method section's reward subsection
(Section~\ref{sec:method-reward}) deliberately forward-references this
discussion. Reward shaping via learned classifiers is increasingly common in
language-model training: deception detectors, motivated-reasoning critics,
style scorers, and reasoning-quality classifiers all enter the reward channel
of one system or another. Our claim is that admitting any such signal as a
Bellman target imports a structural defect that no amount of classifier
quality can fully cure, and that the discipline of keeping reward
verifier-anchored is worth substantial complexity elsewhere.

Consider a learned classifier $f: X \to [0, 1]$ optimized as a reward
signal. After enough updates, the policy converges to behavior that satisfies
$f$ at $f$'s decision boundary; the policy learns the classifier's failure
modes (regions where $f$ disagrees with the underlying property it was
trained on) and exploits them. This is the reward-hacking pattern
characterized formally by \citet{skalse2022reward}, and the proxy-objective
gap it opens has been shown to widen with optimization pressure in the
reward-model overoptimization scaling study of \citet{gao2023scaling}. The
exploitation is not adversarial in intent; it is the natural fixed point of
running an optimizer against an imperfect target. Any imperfection in $f$
becomes an attack surface, and since all learned classifiers are imperfect,
the attack surface is always nonempty. The classical Goodhart pattern
follows: the policy optimizes the proxy at the expense of the underlying
objective, and the gap between proxy and objective is exactly where the
optimizer concentrates its mass.

A verifier $v: X \to \{0, 1\}$ that is not a learned proxy avoids the
worst form of this pathology, but the picture is more nuanced than ``no
Goodhart gap at all'' and deserves a careful distinction. We separate
two cases. \emph{Intrinsic Goodhart} is the failure mode of a learned
proxy: the reward signal is a model fitted to predict the objective,
and under optimization the policy drifts from the objective by
exploiting the proxy's idiosyncrasies. Verifier-anchored signals rule
this case out by construction, because the signal is not a learned
model at all; the SymPy equality call, the unit-test verdict, and the
Lean type-check are deterministic programs, not classifiers, and they
do not have a training distribution off which a policy can wander.
\emph{Extrinsic Goodhart}, sometimes called specification gaming or
the specification-vs-intent gap, is a separate failure mode: the
formally specified signal does not capture the full intended objective,
and optimization exploits the gap. A unit-test harness that only
exercises a narrow slice of input space lets the agent learn solutions
that pass the tests but do not generalize. A SymPy equality check on
canonicalized expressions can be satisfied by an answer that is
mathematically equivalent yet violates a presentation constraint the
problem expects. A Lean kernel accepts any proof term whose types
check, including ones that exploit an unintended axiom or definitional
quirk. In each case the signal and the intended objective are not the
same object, even though the signal is not a learned proxy; the
extrinsic Goodhart channel remains open. Verifier-anchored signals
therefore dominate learned-aesthetic signals on the intrinsic Goodhart
dimension (they close it entirely) and both face the extrinsic Goodhart
dimension. The framework's response to the extrinsic case is to
strengthen the verifier rather than to switch to a learned classifier:
add more test cases, use a formal proof checker rather than a
unit-test harness, tighten canonicalization, audit verifier failures
during training. The agent may still produce outputs the verifier
accepts but a human evaluator would reject, and those cases are
genuine specification failures of the verifier rather than optimization
pathologies; treating verifier coverage as a threat to validity
(Section~\ref{sec:disc-threats}) is the framework's principled way to
keep the gap honest.

The framework therefore reserves Bellman targets for verifier signals.
Learned aesthetic classifiers, when they are useful, enter the system as
MCTS priors (steering early exploration toward trajectories that look
promising on aesthetic grounds) or as expansion-pruning heuristics (blocking
moves that the classifier judges hopelessly off-track). In both roles, the
classifier biases the distribution of trajectories MCTS samples but cannot
rewrite the value function. Bellman backups continue to be computed against
verifier-anchored returns, so the policy that maximizes long-run reward is a
policy that maximizes verifier-anchored success, with aesthetic priors only
shifting which trajectories the agent considered along the way. The asymmetry
is structural: priors influence the search tree's shape; only rewards
determine the converged policy.

We acknowledge that this is a soft commitment rather than a theorem.
Distinguishing verifier-anchored from learned signals in principle is
straightforward; doing it in practice on benchmarks where the verifier
covers only a subset of the desired behavior requires judgment. We invite
empirical testing of whether the discipline pays off: a fair experiment
would train two policies on the same task distribution, one with a learned
critic added to the reward channel and one without, and compare their
behavior on held-out problems where the critic's failure modes are known.
Our prediction is that the verifier-only policy generalizes better off the
critic's training distribution, but the prediction is just that; the
framework here provides the apparatus to run the test, not the test result
itself. The methodological commitment, in any case, is to make the
discipline easy to maintain: by structuring the reward channel so that
inserting a learned critic requires a deliberate choice rather than an
accident, we shift the burden of justification onto whoever would add one.

\subsection{What the Feasibility Probe Surfaced}
\label{sec:disc-pilot}

The preliminary feasibility probe of Section~\ref{sec:pilot} predates the
current action space, so it contributes no benchmark number and we report none
here. Its distinct contribution to the discussion is the forward-looking
monitoring the pre-registered benchmark study of
Section~\ref{sec:experimental-plan} inherits from it, and the demonstration that
the framework's mechanics are implementable end-to-end (a precondition for the
pre-registered study being worth running).

First, composition tracking. A multi-round training-set expansion can change
held-out behavior on families that were not the target of the intervention, and
a composition shift co-occurs with such a change without necessarily being its
cause (Section~\ref{sec:pilot-v3-v4}). Any multi-round result must therefore
report the per-round training composition and the per-class accuracy breakdown
on the held-out set, so that a regression of this shape is attributed to its
actual mechanism rather than to the most visible covariate. The pre-registered
study carries this as a reporting requirement and adds the per-class breakdown
to the plan as an analytic output rather than a post-hoc audit.

Second, target stability. The probe found that visit-count distillation targets
are noisier across collection runs than the distillation prescription assumes
when the search budget is low, with the unseeded LLM transition sampling as the
residual source (Section~\ref{sec:pilot-bsweep}). It also found that the natural
distance metric, the mean $L_1$ between two visit-count distributions, is
confounded by sharpness and can anti-correlate with true stability, so the
correct diagnostic is pairwise recommended-verb argmax agreement across
independent searches from a fixed state. We commit to reporting that diagnostic
on a stratified sample of training states for each benchmark, stratified by
problem family, and to running additional MCTS seeds at the heavy-search budget
on any state where agreement falls below a pre-committed threshold.
\textcolor{red}{\textbf{[PRE-REG TODO: commit the target-stability
threshold (the same number used in Section~\ref{sec:pilot-lessons});
suggested: pairwise recommended-verb agreement $\geq 0.70$ at the
collection budget, stratified by family, confirm]}} The pre-registered
$N_{\mathrm{large}} = 500$ is substantially above the probe's budget and we
expect target stability to improve accordingly, but the expectation is now a
falsifiable monitoring claim rather than an unstated assumption, and the budget
that buys a stable target is set by the slowest-converging family rather than
the average.

The closed-loop exercise surfaced two further methodological commitments the
pre-registered study inherits. The first is honest per-decision compute
accounting: once the frozen model runs in process and one state prefill serves
both the heads and verb generation (Section~\ref{sec:pilot-closedloop}), the
decoder forwards become countable, and the probe found that generation, not
prefill, dominates a decision's token cost. The compute axis of the plan
(Section~\ref{sec:exp-compute}) stays denominated in total inference FLOPs, and
this measurement says where those FLOPs go: the accounting must charge every
generated token inside the search, because a simulation's cost is itself
dominated by generation. The second is that observability and resumability are
research infrastructure, not conveniences: the central claim is an
across-iteration trend (the approximate-$Q$-learning $L_Q$ decrease and the
expert-iteration sharpening of $\ptheta$), and a long run that persists only a
per-iteration $L_{\mathrm{total}}$ summary (in which $L_{\pi}$ near its entropy
floor masks the small $L_Q$ term) cannot surface it. A run long enough to matter
must persist the separated learning signals per iteration and checkpoint often
enough to resume from a crash.

\subsection{Threats to the Empirical Story}
\label{sec:disc-threats}

The experimental plan's threats-to-validity subsection
(Section~\ref{sec:exp-threats}) enumerates the operational concerns; the
discussion here positions them as items to monitor during implementation
and as caveats to any eventual claim. Compute leaks are the most likely
source of an inadvertently inflated result: counting only the LLM forward
passes that occur at root-level decisions, while omitting the forwards
consumed by MCTS rollouts or by chain-of-thought-style verb-instantiator
continuations, would make our method look cheaper than it is at a given
accuracy. The headline frontier figure is meaningful only if every LLM
forward executed during a query is charged against the inference budget
on the $x$-axis. Verifier coverage is the second concern: a lenient or
buggy verifier inflates scores symmetrically across all methods, but it
also lets the agent learn to satisfy the lenient bar rather than the
strict one, so improvements measured under a lenient verifier may not
transfer when the verifier is later tightened.

Benchmark contamination, best-of-$N$ selector choice, and seed variance
round out the list. Contamination is a base-model concern: if the
pretraining corpus of any baseline included the test problems, the
comparison is no longer about deliberation. Selector choice in best-of-$N$
matters because the result depends substantially on whether the selector
is the $Q$-head, a separate learned verifier, or an oracle; we pre-register
the selector choice as part of the baseline grid. Seed variance is the
honest reporting concern: a single-seed headline figure is suggestive at
best, and the plan calls for at least three training seeds with mean and
standard deviation reported. We note these as items the implementation
will need to control, not as defects of the framework.

\subsection{Connection to LLM Deliberation Beyond Math}
\label{sec:disc-llm-deliberation}

Evaluating on verifiable math is a choice of where reward is cleanest, not a
statement that the framework is limited to math. The same structural
commitments transfer to any domain where a non-LLM program supplies the
terminal verdict. Code generation fits naturally: the verifier is a unit-test
suite, the action space gains a \textsc{RunTests} verb playing the role of
\textsc{Verify}, and the terminal reward is the suite's binary outcome.
Formal proof fits even more cleanly: the verifier is the Lean (or Coq, or
Isabelle) kernel, \textsc{Verify} becomes a tactic submission, and the
reward is kernel acceptance. Boolean satisfiability and related decision
procedures admit Z3 (or another SMT solver) as a verifier, the verbs adapting
while the structural pattern is preserved. Restricted fact verification (claims
checkable against a retrieval corpus plus an entailment model whose judgments
themselves type-check structurally) also fits, though the entailment model's
status as verifier rather than learned classifier deserves case-by-case
scrutiny.

What does not fit is open-ended dialogue, creative writing, or any deliberation
task whose terminal reward must come from a learned critic. The framework
accommodates such tasks only after a reward-channel decomposition that admits a
non-learned signal alongside the learned one
(Section~\ref{sec:disc-extensions}). The non-learned signal need not cover
every dimension of quality; it need only be present and be the channel that
participates in Bellman backups. A creative-writing task with a non-learned
signal for ``contains no factual contradictions checkable by external
reference'' could in principle be in scope, with the learned aesthetic critic
supplying MCTS priors but not reward. We do not pursue this here; we flag it as
a direction in which the verifier-anchored reward discipline can travel without
losing its structural advantage.
